\documentclass{article}
\usepackage{hyperref}
\usepackage{Style}

\nocite{*} % Comentar si quiero citar
%\addbibresource{bibliografia.bib} % Quitar el comentado si quiero usar bibliografia

\begin{document}

\begin{minipage}{2.5cm}
    \includegraphics[width=2cm]{imagen_puc.jpg}
\end{minipage}
\begin{minipage}{12cm}
    {\sc Pontificia Universidad Católica de Chile\\
    Facultad de Matemáticas\\
    Departamento de Matemática\\
    Profesor: Manuel Cabezas -- Ayudante: Benjamín Mateluna}
\end{minipage}
\vspace{2ex}

{\centerline{\bf Introducción a la Geometría - MAT1304}
\centerline{\bf Ayudantía 22}}
\centerline{\bf 09 de junio de 2026}

\vspace{1cm}
\noindent\textbf{Problema 1.} Considere en el plano dos circunferencias de centros $A$ y $B$, las
cuales se intersectan en los puntos $C$ y $D$. Demostrar usando solo geometría cartesiana que las
rectas $\overleftrightarrow{AB}$ y $\overleftrightarrow{CD}$ son perpendiculares.

\vspace{5mm}
\noindent\textbf{Problema 2.} Asuma que todas las ecuaciones a continuación representan cónicas. 
Para cada una de ellas encuentre las coordenadas del vértice, centro, focos y ejes:
\begin{enumerate}
    \item $4x^{2}+48y+12x=159$
    \item $x^{2}+4y^{2}-10x-40y+109=0$
    \item $9x^{2}-4y^{2}+54x+16y+29=0$
\end{enumerate}

\vspace{5mm}
\noindent\textbf{Problema 3.} Encontrar una cónica que pase por los puntos $(1,1), (2,2), 
(3,3), (-1,0)$ y $(-2,1)$.

\vspace{5mm}
\noindent\textbf{Problema 4.} Determine todas las rectas tangentes a la circunferencia
\begin{equation*}
    x^{2}-4x+y^{2}-4y+6=0
\end{equation*}
las cuales son paralelas a la recta $x+y=1$.

\vspace{1cm}
\noindent\textbf{\underline{Ejercicios Propuestos:}}

\vspace{5mm}
\noindent\textbf{Extra 1.} Sean $A=(1,0)$ y $B=(2,3)$. Encontrar todos los puntos $C$ tales que
$\triangle ABC$ es equilátero.

\vspace{5mm}
\noindent\textbf{Extra 2.} Encuentre la recta tangente a la parábola $x^{2}=y$ que es paralela
a la recta $y=x+1$.

%\printbibliography % Quitar el comentado si quiero usar bibliografia

\end{document}
